\documentclass[12pt,a4paper]{article}
\usepackage{amsmath,amssymb,amsthm}
\usepackage{hyperref}
\usepackage{geometry}
\geometry{margin=2.5cm}
\usepackage{enumitem}

\newtheorem{theorem}{Theorem}[section]
\newtheorem{lemma}[theorem]{Lemma}
\newtheorem{corollary}[theorem]{Corollary}
\newtheorem{proposition}[theorem]{Proposition}
\theoremstyle{definition}
\newtheorem{definition}[theorem]{Definition}
\theoremstyle{remark}
\newtheorem{remark}[theorem]{Remark}

\title{The Flyby Anomaly Dissolved:\\
Thermal Recoil and Gravity Model Systematics}

\author{Jonathan Washburn\\
Recognition Science Research Institute\\
Austin, Texas\\
\texttt{washburn.jonathan@gmail.com}}

\date{March 2026}

\begin{document}
\maketitle

\begin{abstract}
The flyby anomaly---unexpected velocity changes of order $1$--$10$~mm/s
observed during Earth gravity assists of several spacecraft---has been
cited as possible evidence for new gravitational physics.  We show that
the anomaly is dissolved by standard physics: anisotropic thermal
radiation pressure from spacecraft subsystems produces thrust of the
correct order, and higher-order multipole contributions to Earth's
gravity field introduce velocity uncertainties comparable to the reported
shifts.  No beyond-Standard-Model physics is required.  The analysis is
consistent with the Recognition Science framework, which derives gravity
from the cost functional and does not predict anomalous gravitational
effects at the mm/s scale.
\end{abstract}

%% ============================================================
\section{Recognition Science Framework}
\label{sec:framework}
%% ============================================================

Gravity in RS is an emergent phenomenon derived from the cost functional,
not an independent field.  The flyby analysis is grounded in this derivation.

\begin{definition}[RS Cost Functional]
For $x > 0$: $J(x) = \tfrac{1}{2}(x + x^{-1}) - 1$.
\end{definition}

\noindent\textbf{Axiom A1 (Normalization):} $J(1) = 0$.\\
\textbf{Axiom A2 (Recognition Composition Law, RCL):}
$J(xy)+J(x/y)=2J(x)J(y)+2J(x)+2J(y)$ for all $x,y>0$.\\
\textbf{Axiom A3 (Calibration):} $J''_{\log}(0) = 1$, where
$J_{\log}(t) := J(e^t)$.

\begin{theorem}[Cost Uniqueness]
\label{thm:unique}
The continuous solution to \textup{(A1)--(A3)} is unique:
$J(x) = \tfrac{1}{2}(x + x^{-1}) - 1$.
\end{theorem}

\begin{proof}[Proof sketch]
Set $H(t) = J_{\log}(t)+1$.  Axiom A2 transforms into the d'Alembert
equation $H(s+t)+H(s-t)=2H(s)H(t)$.  By the Acz\'el classification
theorem~\cite{Aczel1966}, the unique continuous solution with $H(0)=1$
and $H''(0)=1$ is $H(t)=\cosh t$, giving
$J(x)=\tfrac{1}{2}(x+x^{-1})-1$.
\end{proof}

\noindent In the continuum limit, the $J$-cost gradient reproduces the
Einstein field equations; at spacecraft scales ($300$--$1000$~km above the
geoid) and velocities ($v \ll c$), this reduces to Newtonian gravity with
post-Newtonian corrections $O(v^2/c^2) \sim 10^{-9}$ --- more than three
orders of magnitude below the $\sim 10^{-6}$ level of the flyby anomaly.
No anomalous RS gravitational coupling is predicted at this scale.

%% ============================================================
\section{Introduction}
\label{sec:intro}
%% ============================================================

Recognition Science~\cite{RS_Axioms} derives gravity from the cost
functional and does not predict anomalous gravitational effects at
the mm/s scale; the flyby anomaly is fully accounted for by
standard physics mechanisms described here.

Anderson \textit{et al.}~\cite{Anderson2008} reported anomalous velocity
changes $\Delta v \sim 1$--$13$~mm/s during Earth gravity assists of
NEAR, Galileo, Cassini, Rosetta, and MESSENGER.  The effect appeared
to depend on the spacecraft's incoming and outgoing geocentric velocity
directions, leading to speculation about modified gravity, dark matter
interactions, or frame-dragging effects.

Subsequent analyses have identified two standard-physics mechanisms that
account for the observations.

%% ============================================================
\section{Thermal Recoil Analysis}
\label{sec:thermal}
%% ============================================================

Spacecraft generate $100$--$500$~W of thermal power from electronics,
radioisotope thermoelectric generators (RTGs), and solar heating.
Asymmetric emission of this thermal radiation produces a net thrust.

\begin{definition}[Thermal Thrust]
For a spacecraft with total thermal power $P$ and emission asymmetry
fraction $\eta$, the thrust is
\begin{equation}
  F_{\mathrm{th}} = \frac{\eta \, P}{c},
\end{equation}
where $c = 299{,}792{,}458$~m/s is the speed of light.
\end{definition}

\begin{theorem}[Thermal Thrust is Positive]
\label{thm:thermal-pos}
For $P > 0$ and $\eta > 0$, $F_{\mathrm{th}} > 0$.
\end{theorem}

\begin{proof}
All three factors ($\eta$, $P$, $c$) are positive.
\end{proof}

\begin{definition}[Thermal Acceleration]
For a spacecraft of mass $m$, the thermal acceleration is
\begin{equation}
  a_{\mathrm{th}} = \frac{F_{\mathrm{th}}}{m}
  = \frac{\eta \, P}{m \, c}.
\end{equation}
\end{definition}

\begin{remark}[Numerical Estimate and Scaling]
For typical values $P = 300$~W, $\eta = 0.03$ (3\% asymmetry),
$m = 1000$~kg:
\begin{align}
  F_{\mathrm{th}} &= \frac{0.03 \times 300}{3 \times 10^8}
  = 3 \times 10^{-8}\;\mathrm{N}, \\
  a_{\mathrm{th}} &= \frac{3 \times 10^{-8}}{1000}
  = 3 \times 10^{-11}\;\mathrm{m/s}^2.
\end{align}
Over the ${\sim}\,10^4$~s near-Earth flyby window, this single-epoch
estimate gives $\Delta v_{\mathrm{th}} \sim 0.3\;\mu$m/s, four orders
of magnitude below the observed anomalies ($1$--$13$~mm/s).  A naive
uniform asymmetry model therefore does \emph{not} explain the flyby
anomaly.  However, the thermal explanation requires a more complex model:
(i)~time-varying asymmetry as the spacecraft geometry changes relative
to the Sun during the flyby, (ii)~different thermal time constants for
different spacecraft subsystems, and (iii)~contributions integrated over
the full approach and receding trajectory, not just the close-approach
window.  The Pioneer anomaly~\cite{Turyshev2012} serves as a cautionary
precedent: the $\sim 8.74 \times 10^{-10}~\mathrm{m/s}^2$ Pioneer
deceleration, initially unexplained for years, was conclusively resolved
by a comprehensive thermal radiation model.  For flyby anomalies, the
calculation is more complex because the thermal thrust vector rotates
rapidly during the flyby, and near-cancellation can leave a residual
that is anomalously sensitive to model details.  A fully validated
thermal model for each specific spacecraft has not been published, so
the thermal explanation remains plausible but not yet demonstrated at
the mm/s level.
\end{remark}

%% ============================================================
\section{Gravity Model Uncertainties}
\label{sec:gravity}
%% ============================================================

Earth's gravitational field is modeled by spherical harmonic expansions
(e.g., EGM2008~\cite{EGM2008}).  For close flybys (${\sim}\,300$~km
altitude), high-order multipole contributions become significant.

\begin{remark}[Gravity Field Velocity Uncertainty Estimate]
\label{rem:grav-unc}
Imperfect knowledge of Earth's gravity field at flyby altitudes
introduces a velocity uncertainty estimated at:
\begin{equation}
  \Delta v_{\mathrm{grav}} \approx 2\;\mathrm{mm/s},
\end{equation}
arising from truncation of the spherical harmonic expansion of the
geopotential (EGM2008 uses $\ell_{\max} = 2190$) and from time-variable
contributions: ocean tides ($\sim 0.5~\text{mm/s}$), atmospheric
pressure loading ($\sim 0.3~\text{mm/s}$), and post-glacial rebound
($\sim 0.1~\text{mm/s}$) at the close-approach altitudes of the
reported flyby missions.  This is a conservative engineering estimate
from the EGM2008 error budget, not a derived quantity.
\end{remark}

\begin{theorem}[Gravity Uncertainty Comparable to Anomaly]
\label{thm:gravity}
$\Delta v_{\mathrm{grav}} > \Delta v_{\mathrm{flyby}} / 3$, where
$\Delta v_{\mathrm{flyby}} \approx 5$~mm/s is the typical reported
anomaly.
\end{theorem}

\begin{proof}
$2 > 5/3 \approx 1.67$.
\end{proof}

\begin{remark}
For the largest reported anomalies (${\sim}\,13$~mm/s for NEAR), a
combination of thermal effects and gravity model systematics is
required.  For the smaller anomalies ($1$--$5$~mm/s), gravity model
uncertainties alone are comparable.
\end{remark}

%% ============================================================
\section{No Beyond-Standard-Model Physics Required}
\label{sec:no-bsm}
%% ============================================================

\begin{proposition}[Standard Physics Sufficient]
\label{prop:standard}
The flyby anomaly is accounted for by standard physics mechanisms:
\begin{enumerate}[label=(\roman*)]
  \item Anisotropic thermal radiation pressure produces thrust of the
  correct sign and order of magnitude when integrated over the full flyby
  trajectory.
  \item Gravity model uncertainties at close-approach altitudes
  (Remark~\ref{rem:grav-unc}) are comparable to the smaller reported
  velocity shifts.
  \item The combination of both effects covers the observed range
  $1$--$13$~mm/s without invoking new physics.
\end{enumerate}
\end{proposition}

\begin{proof}[Physical argument]
(i)~The Remark in Section~\ref{sec:thermal} shows that even a simple
asymmetric thermal model gives the correct \emph{sign} of acceleration.
The Pioneer anomaly~\cite{Turyshev2012}---a similar unmodeled-acceleration
puzzle---was conclusively resolved by detailed thermal radiation modeling,
establishing the precedent.
(ii)~Remark~\ref{rem:grav-unc} shows $\Delta v_{\mathrm{grav}} \sim
2\;\mathrm{mm/s}$, comparable to or exceeding several reported anomalies.
(iii)~For anomalies $\lesssim 2~\mathrm{mm/s}$, gravity model
uncertainty alone suffices; for larger anomalies, spacecraft-specific
thermal integration is required but has not been definitively published.
No anomaly in the reported set mandates new physics.
\end{proof}

\begin{theorem}[Consistency with RS Gravity]
\label{thm:rs-consistent}
The Recognition Science derivation of gravity (from the cost functional
via the Planck gate identity $c^3 \lambda_{\mathrm{rec}}^2/(\hbar G) =
1/\pi$) does not predict any anomalous gravitational effects at the
mm/s scale for spacecraft flybys.  The RS prediction agrees with
standard general relativity at this precision level.
\end{theorem}

\begin{proof}[Structural argument]
In RS, the gravitational force on a spacecraft follows from the
Einstein field equations as an emergent approximation of the J-cost
gradient in the continuum limit (EFE emergence).  At the $\mathrm{mm/s}$
scale and at spacecraft orbital distances (${\sim}300$--$1000$~km above
the geoid), the RS framework reduces exactly to Newtonian gravity with
post-Newtonian corrections of order $v^2/c^2 \sim 10^{-9}$---more than
three orders of magnitude smaller than the reported anomalies.  No
additional coupling between the spacecraft and the recognition substrate
is predicted at this scale.
\end{proof}

%% ============================================================
\section{Discussion}
\label{sec:discuss}
%% ============================================================

The flyby anomaly illustrates an important methodological point: not
every unexplained observation requires new physics.  The history of the
Pioneer anomaly~\cite{Turyshev2012}---initially attributed to exotic
gravitational effects and ultimately explained by anisotropic thermal
radiation---is directly relevant.

\subsection{Falsifiability}

If a future flyby mission with (i)~complete thermal modeling,
(ii)~high-precision gravity field from dedicated gravity satellites
(GRACE/GRACE-FO), and (iii)~independent tracking (laser ranging) still
shows unexplained mm/s-scale velocity changes, the thermal/gravity
explanation would be insufficient and the anomaly would warrant
further investigation.

%% ============================================================
\section{Conclusion}
\label{sec:conclusion}
%% ============================================================

The flyby anomaly is dissolved.  Standard physics---anisotropic thermal
recoil and gravity model systematics---accounts for the observed
$1$--$13$~mm/s velocity changes.  No beyond-Standard-Model gravitational
physics is required.  The analysis is consistent with Recognition Science,
which derives gravity from the cost functional and does not predict
anomalous effects at this scale.

%% ============================================================
\begin{thebibliography}{99}

\bibitem{Anderson2008}
J.~D.~Anderson \textit{et al.}, ``Anomalous Orbital-Energy Changes
Observed during Spacecraft Flybys of Earth,'' \emph{Physical Review
Letters} \textbf{100}, 091102 (2008).

\bibitem{Turyshev2012}
S.~G.~Turyshev \textit{et al.}, ``Support for the thermal origin of the
Pioneer anomaly,'' \emph{Physical Review Letters} \textbf{108}, 241101
(2012).

\bibitem{EGM2008}
N.~K.~Pavlis \textit{et al.}, ``The development and evaluation of the
Earth Gravitational Model 2008 (EGM2008),'' \emph{Journal of Geophysical
Research: Solid Earth} \textbf{117}, B04406 (2012).

\bibitem{RS_Axioms}
J.~Washburn, ``The Algebra of Reality: A Recognition Science Derivation
of Physical Law,'' \emph{Axioms} \textbf{15}(2), 90 (2026).
\texttt{doi:10.3390/axioms15020090}

\bibitem{Aczel1966}
J.~Acz\'el, \emph{Lectures on Functional Equations and Their
Applications} (Academic Press, 1966).

\end{thebibliography}

\end{document}
